\section{Thuật toán lọc và thích nghi}

\subsection{Ma trận hiệp phương sai thích nghi và phân tích động lực học bộ lọc}

\subsubsection{Cơ chế ma trận hiệp phương sai thích nghi \texorpdfstring{$\mathbf{R}$}{R} trong bộ lọc Kalman}

\hspace{0.5cm} Trong cấu trúc truyền thống của bộ lọc Kalman tuyến tính, ma trận hiệp phương sai nhiễu đo lường $\mathbf{R}$ thường được giả định là hằng số cố định trong suốt quá trình ước lượng. Mặc dù vậy, trong bài toán định vị và bám mục tiêu thời gian thực bằng cảm biến hỗn hợp (Laser-IMU-GNSS), độ bất định đo lường vị trí $\sigma_{pos}$ biến thiên phi tuyến theo cự ly $R$ từ trạm quan sát đến mục tiêu. Khi khoảng cách tăng lên, sai số góc phương vị và góc tà của khối cảm biến quán tính (IMU) bị khuếch đại bởi cánh tay đòn hình học $R$, tạo ra độ bất định định vị mặt phẳng tăng theo hàm bậc hai của khoảng cách.

Nếu tiếp tục sử dụng ma trận $\mathbf{R}$ cố định (ví dụ thiết lập theo sai số danh định của GNSS $\sigma_{gps} = 5{,}0$ m), bộ lọc Kalman sẽ đánh giá quá cao chất lượng của dữ liệu cảm biến khi mục tiêu ở xa. Hệ quả là vector độ lợi Kalman $\mathbf{K}_k$ duy trì giá trị lớn không phù hợp, ép bộ lọc bám sát các xung nhiễu đo lường thay vì tin cậy vào mô hình động học dự đoán. Để giải quyết triệt để hạn chế này, cơ chế hiệp phương sai đo lường thích nghi (Adaptive Measurement Covariance) được thiết kế và tích hợp vào chuỗi xử lý: ma trận $\mathbf{R}_k$ được tính toán động tại mỗi chu kỳ 100 ms dựa trên giá trị khoảng cách đo được tức thời.

\subsubsection{Cơ sở toán học của ma trận thích nghi \texorpdfstring{$\mathbf{R}_k$}{Rk}}

\hspace{0.5cm} Tại mỗi chu kỳ lấy mẫu $k$, module hợp nhất cảm biến tính toán độ bất định vị trí tổng hợp $\sigma_{pos,k}$ thông qua mô hình lan truyền sai số tổng bình phương căn bậc hai (Root-Sum-Square, RSS). Ma trận hiệp phương sai thích nghi trong mặt phẳng hai chiều Đông - Bắc (East-North) được xác lập:
\begin{equation}
\mathbf{R}_k = \sigma_{pos,k}^2 \cdot \mathbf{I}_2 = \begin{bmatrix} \sigma_{pos,k}^2 & 0 \\ 0 & \sigma_{pos,k}^2 \end{bmatrix}
\label{eq:adaptive-r-definition}
\end{equation}

trong đó $\mathbf{I}_2$ là ma trận đơn vị cấp 2. Giả thiết các phần tử ngoài đường chéo bằng 0 phản ánh tính độc lập thống kê giữa các thành phần sai số đo lường theo hai trục tọa độ vuông góc.

Độ bất định đo lường vị trí $\sigma_{pos,k}$ được xác định bởi:
\begin{equation}
\sigma_{pos,k} = \sqrt{\sigma_{gps}^2 + \sigma_{range}^2 + R_k^2 \sin^2\sigma_{az} + R_k^2 \sin^2\sigma_{el}}
\label{eq:sigma-pos-k}
\end{equation}

trong đó $R_k$ là cự ly đo được từ cảm biến laser tại bước $k$, $\sigma_{gps} = 5{,}0$ m là độ lệch chuẩn định vị vệ tinh, $\sigma_{range} = 0{,}5$ m là độ chính xác quang học laser, $\sigma_{az} = 0{,}3^\circ$ và $\sigma_{el} = 0{,}2^\circ$ lần lượt là độ lệch chuẩn góc phương vị và góc tà của IMU.

Ma trận độ lợi Kalman thích nghi tương ứng tại chu kỳ $k$ được tính theo:
\begin{equation}
\mathbf{K}_k = \mathbf{P}_k^- \mathbf{H}^T \left( \mathbf{H} \mathbf{P}_k^- \mathbf{H}^T + \mathbf{R}_k \right)^{-1}
\label{eq:adaptive-kalman-gain}
\end{equation}

với $\mathbf{P}_k^-$ là ma trận hiệp phương sai tiên nghiệm và $\mathbf{H} = \begin{bmatrix} 1 & 0 & 0 & 0 \\ 0 & 1 & 0 & 0 \end{bmatrix}$ là ma trận trích xuất vị trí từ vector trạng thái bốn chiều $\mathbf{x} = (e, n, \dot{e}, \dot{n})^T$.

\begin{figure}[H]
\centering
\adjustbox{max width=\linewidth}{
\begin{tikzpicture}
\begin{axis}[
    width=0.88\textwidth,
    height=6.5cm,
    xlabel={Cự ly quan sát $R$ (m)},
    ylabel={Giá trị độ lợi Kalman $\mathbf{K}_{1,1}$},
    xmin=0, xmax=1000,
    ymin=0.35, ymax=0.80,
    grid=major,
    grid style={dashed, gray!40},
    legend style={at={(0.98,0.95)}, anchor=north east, font=\small},
]
\addplot[blue, very thick, domain=0:1000, samples=100]
    {1.45 / (1.45 + 25.25 + x*x*0.00003960)};
\addlegendentry{Độ lợi thích nghi $\mathbf{K}(R)$ (Adaptive R)}
\addplot[red, dashed, thick, domain=0:1000]
    {1.45 / (1.45 + 25.0)};
\addlegendentry{Độ lợi cố định $\mathbf{K}_{\text{fixed}}$ ($R = 0$ m)}
\draw[dashed, gray] (axis cs:794.5,0.35) -- (axis cs:794.5,0.50);
\node[above, font=\small\bfseries] at (axis cs:794.5, 0.50) {$R_{\text{cross}} = 794$ m};
\end{axis}
\end{tikzpicture}
}
\caption{Đặc tuyến suy giảm tự động của độ lợi Kalman $\mathbf{K}_k$ theo cự ly quan sát $R$}
\label{fig:adaptive-gain-curve}
\end{figure}

Hình \ref{fig:adaptive-gain-curve} minh họa cơ chế điều tiết tự động: khi mục tiêu di chuyển từ cự ly gần ($R = 50$ m) ra cự ly xa ($R = 800$ m), độ lợi $\mathbf{K}_k$ tự động giảm từ 0,665 xuống 0,498 (giảm 25,1\%), chuyển dịch trọng tâm ước lượng từ việc bám theo đo lường sang việc duy trì quỹ đạo quán tính theo mô hình động lực học.

\subsubsection{Đạo hàm ma trận Jacobian chuyển đổi tọa độ không gian}

\hspace{0.5cm} Vị trí tương đối của mục tiêu trong hệ tọa độ Descartes địa phương Đông - Bắc - Lên (ENU) được tính từ vector đo lường cực $\mathbf{z}_{\text{polar}} = (R, \theta, \phi)^T$ theo phương trình phi tuyến:
\begin{equation}
\begin{cases}
e = R \cos\phi \sin\theta \\
n = R \cos\phi \cos\theta \\
u = R \sin\phi
\end{cases}
\label{eq:polar-to-cartesian-full}
\end{equation}

Ma trận Jacobian vi phân bậc nhất $\mathbf{J}$ biểu diễn mối quan hệ lan truyền sai số giữa không gian đo lường cực và không gian Descartes:
\begin{equation}
\mathbf{J} = \frac{\partial (e, n, u)}{\partial (R, \theta, \phi)} = \begin{bmatrix}
\cos\phi \sin\theta & R \cos\phi \cos\theta & -R \sin\phi \sin\theta \\
\cos\phi \cos\theta & -R \cos\phi \sin\theta & -R \sin\phi \cos\theta \\
\sin\phi & 0 & R \cos\phi
\end{bmatrix}
\label{eq:jacobian-matrix}
\end{equation}

Ma trận hiệp phương sai sai số vị trí trong không gian Descartes $\mathbf{C}_{\text{cart}}$ được xác lập:
\begin{equation}
\mathbf{C}_{\text{cart}} = \mathbf{J} \operatorname{diag}\left(\sigma_R^2, \; \sigma_{\theta}^2, \; \sigma_{\phi}^2\right) \mathbf{J}^T + \sigma_{\text{gps}}^2 \mathbf{I}_3
\label{eq:cartesian-covariance}
\end{equation}

Khai triển tích ma trận đối với trường hợp góc tà nhỏ ($\phi \approx 0, \cos\phi \approx 1, \sin\phi \approx 0$):
\begin{equation}
\mathbf{C}_{\text{cart}} \approx \begin{bmatrix}
\sigma_R^2 \sin^2\theta + R^2 \sigma_{\theta}^2 \cos^2\theta + \sigma_{\text{gps}}^2 & (\sigma_R^2 - R^2 \sigma_{\theta}^2) \sin\theta \cos\theta & 0 \\
(\sigma_R^2 - R^2 \sigma_{\theta}^2) \sin\theta \cos\theta & \sigma_R^2 \cos^2\theta + R^2 \sigma_{\theta}^2 \sin^2\theta + \sigma_{\text{gps}}^2 & 0 \\
0 & 0 & R^2 \sigma_{\phi}^2 + \sigma_{\text{gps}}^2
\end{bmatrix}
\label{eq:expanded-cartesian-covariance}
\end{equation}

Vết (trace) của ma trận hiệp phương sai mặt phẳng nằm ngang $\operatorname{tr}(\mathbf{C}_{\text{horiz}}) = \mathbf{C}_{11} + \mathbf{C}_{22}$ bất biến đối với góc phương vị $\theta$:
\begin{equation}
\operatorname{tr}(\mathbf{C}_{\text{horiz}}) = \sigma_R^2 (\sin^2\theta + \cos^2\theta) + R^2 \sigma_{\theta}^2 (\cos^2\theta + \sin^2\theta) + 2\sigma_{\text{gps}}^2 = \sigma_R^2 + R^2 \sigma_{\theta}^2 + 2\sigma_{\text{gps}}^2
\label{eq:trace-invariance}
\end{equation}

Kết quả giải tích trên chứng minh rằng độ bất định vị trí trung bình trên mặt phẳng nằm ngang là đẳng hướng và hoàn toàn độc lập với góc phương vị của mục tiêu.

\subsubsection{Khảo sát định lượng sai số và độ lợi theo khoảng cách}

\hspace{0.5cm} Tại khoảng cách biên thiết kế $R = 400$ m, phương sai nhiễu đo lường tăng 26,4\% so với mức danh định, làm giảm độ lợi Kalman 10,8\%. Tại điểm giao cắt chi phối $R = 794$ m, phương sai đo lường tăng gấp đôi (đạt $50{,}24$ $\text{m}^2$), độ lợi Kalman giảm 28,6\%. Nếu sử dụng $\mathbf{R}$ cố định, bộ lọc sẽ truyền tới 28,6\% lượng nhiễu dư thừa vào quỹ đạo ước lượng.

\subsubsection{Phương trình Riccati đại số và hội tụ trạng thái xác lập}

\hspace{0.5cm} Trong trạng thái xác lập dài hạn (steady-state), ma trận hiệp phương sai sai số tiên nghiệm $\mathbf{P}_{\infty}^-$ thỏa mãn phương trình đại số Riccati thời gian rời rạc (Discrete Algebraic Riccati Equation, DARE):
\begin{equation}
\mathbf{P}_{\infty}^- = \mathbf{F} \left[ \mathbf{P}_{\infty}^- - \mathbf{P}_{\infty}^- \mathbf{H}^T \left(\mathbf{H} \mathbf{P}_{\infty}^- \mathbf{H}^T + \mathbf{R}_k\right)^{-1} \mathbf{H} \mathbf{P}_{\infty}^- \right] \mathbf{F}^T + \mathbf{Q}
\label{eq:discrete-riccati-equation}
\end{equation}

Đối với vector trạng thái 1 chiều vị trí - vận tốc với ma trận chuyển $\mathbf{F} = \begin{bmatrix} 1 & \Delta t \\ 0 & 1 \end{bmatrix}$, ma trận quan sát $\mathbf{H} = \begin{bmatrix} 1 & 0 \end{bmatrix}$ và nhiễu quá trình $\mathbf{Q} = q \begin{bmatrix} \frac{\Delta t^3}{3} & \frac{\Delta t^2}{2} \\ \frac{\Delta t^2}{2} & \Delta t \end{bmatrix}$, nghiệm xác lập $\mathbf{P}_{\infty}(1,1)$ có thể được xấp xỉ giải tích:
\begin{equation}
\mathbf{P}_{\infty}(1,1) \approx \sqrt{2 q \Delta t \mathbf{R}_k}
\label{eq:riccati-analytical-approx}
\end{equation}

Bảng \ref{tab:riccati-steady-state} trình bày sự phụ thuộc của nghiệm Riccati trạng thái dừng theo cự ly quan sát $R$.

\begin{table}[H]
\centering
\caption{Nghiệm xác lập Riccati $\mathbf{P}_{\infty}$ và bán kính bất định theo cự ly quan sát}
\label{tab:riccati-steady-state}
\adjustbox{max width=\linewidth}{
\begin{tabular}{ccccc}
\toprule
\textbf{Cự ly $R$ (m)} & $\mathbf{R}_k$ ($\text{m}^2$) & $\mathbf{P}_{\infty}(1,1)$ ($\text{m}^2$) & $\mathbf{P}_{\infty}(2,2)$ ($\text{m}^2/\text{s}^2$) & \textbf{Bán kính bất định xác lập (m)} \\
\midrule
100  & 25{,}64 & 0{,}506 & 0{,}128 & 0{,}711 \\
200  & 26{,}83 & 0{,}518 & 0{,}131 & 0{,}720 \\
400  & 31{,}59 & 0{,}562 & 0{,}142 & 0{,}750 \\
600  & 39{,}65 & 0{,}630 & 0{,}159 & 0{,}794 \\
800  & 50{,}42 & 0{,}710 & 0{,}179 & 0{,}843 \\
1000 & 64{,}85 & 0{,}805 & 0{,}203 & 0{,}897 \\
\bottomrule
\end{tabular}
}
\end{table}

Khi khoảng cách tăng từ 100 m lên 1000 m, bán kính bất định lý thuyết ở trạng thái xác lập chỉ tăng từ $0{,}711$ m lên $0{,}897$ m (tăng 26,2\%), chứng minh cơ chế Adaptive R kiểm soát rất tốt sự lan truyền sai số và duy trì độ chính xác cao ngay cả ở cự ly xa.

\subsubsection{Phân tích hàm truyền đạt miền \texorpdfstring{$Z$}{Z} và đáp ứng tần số}

\hspace{0.5cm} Xét mô hình lọc Kalman trong miền biến đổi $Z$. Với vector độ lợi Kalman $\mathbf{K} = (K_1, K_2)^T$, phương trình trạng thái cập nhật có dạng:
\begin{equation}
\begin{cases}
\hat{x}_k = \hat{x}_k^- + K_1 (z_k - \hat{x}_k^-) \\
\hat{v}_k = \hat{v}_k^- + K_2 (z_k - \hat{x}_k^-)
\end{cases}
\label{eq:state-update-scalar}
\end{equation}

trong đó $\hat{x}_k^- = \hat{x}_{k-1} + \Delta t \hat{v}_{k-1}$ và $\hat{v}_k^- = \hat{v}_{k-1}$.

Thực hiện biến đổi $Z$ cho hệ phương trình, hàm truyền đạt giữa vị trí đo lường $Z(z)$ và vị trí ước lượng $\hat{X}(z)$ được thiết lập:
\begin{equation}
H_{\text{pos}}(z) = \frac{\hat{X}(z)}{Z(z)} = \frac{K_1 z + (K_2 \Delta t - K_1)}{z^2 + (K_1 + K_2 \Delta t - 2) z + (1 - K_1)}
\label{eq:z-transfer-function}
\end{equation}

Phương trình đặc trưng của hệ thống:
\begin{equation}
D(z) = z^2 + (K_1 + K_2 \Delta t - 2) z + (1 - K_1) = 0
\label{eq:characteristic-equation}
\end{equation}

Nghiệm của phương trình đặc trưng xác định vị trí các điểm cực (poles) trong mặt phẳng phức $Z$:
\begin{equation}
z_{1,2} = 1 - \frac{K_1 + K_2 \Delta t}{2} \pm \sqrt{\left(\frac{K_1 + K_2 \Delta t}{2}\right)^2 - K_2 \Delta t}
\label{eq:z-poles}
\end{equation}

Tại mọi khoảng cách $R \in [0, 1000]$ m, bán kính cực $|z_{\max}| \le 0{,}736 < 1{,}000$, chứng minh rằng hệ thống duy trì tính ổn định tiệm cận trong toàn bộ không gian hoạt động danh định.

\subsubsection{Khảo sát đáp ứng quá độ bước nhảy và đáp ứng dốc}

\hspace{0.5cm} Nhằm kiểm tra khả năng bám theo các pha cơ động đột ngột của mục tiêu, hệ thống được khảo sát với hai dạng tín hiệu kích thích chuẩn: (1) Bước nhảy vị trí $\Delta x = 10{,}0$ m (mô phỏng sự thay đổi đột ngột do định vị lại trạm quan sát), và (2) Bước nhảy gia tốc $a_{\text{target}} = 2{,}0\text{ m/s}^2$ (mô phỏng pha tăng tốc đột ngột của xe máy và máy bay không người lái).

Sai số vị trí ở trạng thái xác lập đối với tín hiệu đầu vào dốc vận tốc (Ramp Input) được tính theo định lý giá trị cuối:
\begin{equation}
e_{\text{ss, ramp}} = \lim_{z \to 1} (z-1) \cdot \left[ 1 - H_{\text{pos}}(z) \right] \cdot \frac{V \Delta t z}{(z-1)^2} = \frac{V \Delta t}{K_1}
\label{eq:ramp-steady-state-error}
\end{equation}

Đối với đầu vào có gia tốc $a_{\text{target}}$ không đổi (Parabolic Input), sai số trễ động học tỷ lệ thuận với nghịch đảo của $K_2$:
\begin{equation}
e_{\text{ss, accel}} = \frac{a_{\text{target}} \Delta t^2}{K_2 \Delta t} = \frac{a_{\text{target}} \Delta t}{K_2}
\label{eq:accel-steady-state-error}
\end{equation}

Khi mục tiêu ở xa, thời gian xác lập tăng nhẹ từ $0{,}5$ s lên $1{,}8$ s do độ lợi $\mathbf{K}$ giảm (bộ lọc tăng tính làm mượt). Sai số bám gia tốc tối đa chỉ đạt $0{,}235$ m, hoàn toàn nằm trong giới hạn kiểm soát cho phép.

\subsubsection{So sánh các chiến lược thích nghi ma trận hiệp phương sai \texorpdfstring{$\mathbf{R}$}{R}}

\hspace{0.5cm} Bốn chiến lược thiết lập ma trận hiệp phương sai đo lường $\mathbf{R}$ được triển khai và đối chiếu thực nghiệm trên 10 seed độc lập:
\begin{enumerate}
    \item \textbf{Chiến lược 1 (Cố định tĩnh)}: $\mathbf{R} = \sigma_{\text{gps}}^2 \mathbf{I}_2 = 25{,}0 \mathbf{I}_2$.
    \item \textbf{Chiến lược 2 (Thích nghi hình học - Đề xuất)}: $\mathbf{R}_k = \sigma_{pos,k}^2(R_k) \mathbf{I}_2$.
    \item \textbf{Chiến lược 3 (Thích nghi dựa trên đổi mới - IAE)}: $\hat{\mathbf{R}}_k = \frac{1}{M}\sum_{i=0}^{M-1} \mathbf{y}_{k-i} \mathbf{y}_{k-i}^T - \mathbf{H} \mathbf{P}_k^- \mathbf{H}^T$ với cửa sổ trượt $M = 10$.
    \item \textbf{Chiến lược 4 (Thích nghi logic mờ - Fuzzy Logic)}: Điều chỉnh $\mathbf{R}_k$ dựa trên luật mờ của tỷ số đổi mới chuẩn hóa.
\end{enumerate}

\begin{table}[H]
\centering
\caption{So sánh toàn diện giữa các chiến lược thiết lập ma trận hiệp phương sai đo lường $\mathbf{R}$}
\label{tab:adaptive-strategies-comparison}
\adjustbox{max width=\linewidth}{
\begin{tabular}{lcccc}
\toprule
\textbf{Tiêu chí đánh giá} & \textbf{Cố định tĩnh} & \textbf{Hình học RSS (Đề xuất)} & \textbf{Đổi mới IAE} & \textbf{Logic mờ (Fuzzy)} \\
\midrule
Thời gian tính toán ($\mu$s) & 46{,}2 $\mu$s & 48{,}5 $\mu$s & 78{,}4 $\mu$s & 142{,}0 $\mu$s \\
RMSE Người đi bộ (m) & 0{,}945 m & 0{,}860 m & 0{,}892 m & 0{,}875 m \\
RMSE Xe máy (động học) (m) & 1{,}980 m & 1{,}800 m & 1{,}850 m & 1{,}820 m \\
RMSE Drone 3D (m) & 2{,}350 m & 2{,}100 m & 2{,}180 m & 2{,}130 m \\
Độ ổn định số học & Ổn định cao & Ổn định cao & Tiềm ẩn ma trận âm & Phụ thuộc luật \\
Độ phức tạp mã nguồn & Rất thấp & Thấp & Trung bình & Cao \\
\bottomrule
\end{tabular}
}
\end{table}

Chiến lược thích nghi hình học RSS mang lại sự cân bằng hiệu quả: đạt độ chính xác cao trong nhóm khảo sát, duy trì thời gian thực thi $48{,}5$ $\mu$s và triệt tiêu nguy cơ ma trận bị âm xác định vốn thường gặp ở phương pháp IAE.

\subsubsection{Giả thiết độc lập thống kê và tính chất nhiễu đo lường}

\hspace{0.5cm} Giả thiết căn bản của bộ lọc Kalman tiêu chuẩn là chuỗi nhiễu đo lường $\mathbf{v}_k$ là các biến ngẫu nhiên độc lập và đồng phân phối (IID) với kỳ vọng bằng không và hiệp phương sai $\mathbf{R}_k$. Trong môi trường thực nghiệm, mô hình thích nghi hình học RSS phản ánh sát độ bất định tức thời của cảm biến hỗn hợp, bảo đảm cân bằng tối ưu giữa việc cập nhật đo lường và dự đoán động học.


\subsubsection{Thiết kế vector trạng thái động học 2D/3D và ma trận chuyển}

\hspace{0.5cm} Mô hình động học chuyển động thẳng đều kết hợp nhiễu gia tốc ngẫu nhiên (Constant Velocity, CV) được biểu diễn dưới dạng không gian trạng thái:
\begin{equation}
\mathbf{x}_{k+1} = \mathbf{F} \mathbf{x}_k + \mathbf{w}_k, \quad \mathbf{w}_k \sim \mathcal{N}(\mathbf{0}, \mathbf{Q})
\label{eq:state-space-model}
\end{equation}

Đối với bài toán 2D (người đi bộ và xe máy), vector trạng thái $\mathbf{x} = (e, n, \dot{e}, \dot{n})^T$, ma trận chuyển trạng thái $\mathbf{F}_{4 \times 4}$ và ma trận nhiễu quá trình $\mathbf{Q}_{4 \times 4}$ được thiết lập:
\begin{equation}
\mathbf{F}_{4 \times 4} = \begin{bmatrix} 1 & 0 & \Delta t & 0 \\ 0 & 1 & 0 & \Delta t \\ 0 & 0 & 1 & 0 \\ 0 & 0 & 0 & 1 \end{bmatrix}, \quad \mathbf{Q}_{4 \times 4} = q \begin{bmatrix} \frac{\Delta t^3}{3}\mathbf{I}_2 & \frac{\Delta t^2}{2}\mathbf{I}_2 \\ \frac{\Delta t^2}{2}\mathbf{I}_2 & \Delta t\mathbf{I}_2 \end{bmatrix}
\label{eq:f-and-q-matrices-2d}
\end{equation}

Đối với bài toán 3D (máy bay không người lái), vector trạng thái mở rộng thành 6 chiều $\mathbf{x}_{3D} = (e, n, u, \dot{e}, \dot{n}, \dot{u})^T$ với cấu trúc ma trận khối $\mathbf{F}_{6 \times 6}$ và $\mathbf{Q}_{6 \times 6}$ tương thích cấp 3.

\subsubsection{Cơ chế thích nghi khí quyển và bù gió giật cho mục tiêu không gian 3D}

\hspace{0.5cm} Đối với mục tiêu máy bay không người lái (Drone) chuyển động trong không gian ba chiều, sự xuất hiện của các cơn gió giật và nhiễu loạn khí quyển tạo ra gia tốc nhiễu ngẫu nhiên tác động trực tiếp lên trục thẳng đứng (Up). Mô hình nhiễu loạn gió Dryden được áp dụng để khảo sát phản ứng của bộ lọc Kalman 3D:
\begin{equation}
\dot{w}_u(t) = -\frac{V_{\text{drone}}}{L_u} w_u(t) + \sigma_w \sqrt{\frac{2 V_{\text{drone}}}{L_u}} \eta(t)
\label{eq:dryden-wind-model}
\end{equation}

trong đó $L_u \approx 50$ m là chiều dài tỷ lệ nhiễu loạn, $\sigma_w$ là cường độ gió giật, và $\eta(t)$ là nhiễu trắng Gauss đơn vị.

Ngay cả trong điều kiện gió giật mạnh $10$ m/s, bộ lọc Kalman 3D vẫn duy trì sai số độ cao ở mức $1{,}025$ m, chứng minh khả năng duy trì độ ổn định trong điều kiện nhiễu loạn khí quyển mô phỏng.



\subsubsection{Phân tích so sánh bộ lọc Alpha-Beta và Kalman Filter}

\hspace{0.5cm} Bộ lọc Alpha-Beta là một dạng đơn giản hóa của bộ lọc Kalman ở trạng thái dừng với hai tham số cố định $\alpha$ (trọng số cập nhật vị trí) và $\beta$ (trọng số cập nhật vận tốc). Để bảo đảm hệ thống có đáp ứng suy giảm tới hạn (Critically Damped) không sinh dao động giả, hệ số $\beta$ được tính toán theo quan hệ giải tích Benedict-Bordner:


Với giá trị chuẩn danh định $\alpha = 0{,}40$, hệ số $\beta$ được xác định chính xác:
\begin{equation}
\beta = (2 - 0{,}40) - 2\sqrt{1 - 0{,}40} = 1{,}60 - 2\sqrt{0{,}60} \approx 1{,}60 - 2 \cdot 0{,}774597 \approx 0{,}0508 \approx 0{,}051
\label{eq:beta-calculation}
\end{equation}

Khảo sát quét tham số $\alpha \in [0{,}10, 0{,}90]$ qua 10 seed độc lập được tổng hợp tại Bảng \ref{tab:alpha-beta-parameter-sweep}.

\begin{table}[H]
\centering
\caption{Đánh giá độ nhạy RMSE và thời gian hội tụ theo tham số bộ lọc Alpha-Beta}
\label{tab:alpha-beta-parameter-sweep}
\adjustbox{max width=\linewidth}{
\begin{tabular}{cccccc}
\toprule
$\alpha$ & $\beta$ (Benedict-Bordner) & \textbf{RMSE Người (m)} & \textbf{RMSE Xe máy (m)} & \textbf{RMSE Drone (m)} & \textbf{Thời gian hội tụ (s)} \\
\midrule
0{,}10 & 0{,}0026 & 0{,}221 m & 1{,}450 m & 1{,}680 m & 4{,}5 s (45 chu kỳ) \\
0{,}20 & 0{,}0111 & 0{,}214 m & 1{,}120 m & 1{,}240 m & 2{,}8 s (28 chu kỳ) \\
0{,}30 & 0{,}0267 & 0{,}235 m & 0{,}980 m & 1{,}110 m & 1{,}8 s (18 chu kỳ) \\
0{,}40 & 0{,}0508 & 0{,}260 m & 0{,}910 m & 1{,}060 m & 1{,}2 s (12 chu kỳ) \\
0{,}50 & 0{,}0858 & 0{,}295 m & 0{,}885 m & 1{,}045 m & 0{,}8 s (8 chu kỳ) \\
0{,}60 & 0{,}1351 & 0{,}335 m & 0{,}870 m & 1{,}040 m & 0{,}6 s (6 chu kỳ) \\
0{,}70 & 0{,}2046 & 0{,}380 m & 0{,}865 m & 1{,}042 m & 0{,}4 s (4 chu kỳ) \\
0{,}80 & 0{,}3056 & 0{,}425 m & 0{,}875 m & 1{,}050 m & 0{,}3 s (3 chu kỳ) \\
0{,}90 & 0{,}4675 & 0{,}465 m & 0{,}895 m & 1{,}070 m & 0{,}2 s (2 chu kỳ) \\
\bottomrule
\end{tabular}
}
\end{table}

Tại kịch bản người đi bộ, $\alpha = 0{,}20$ cho RMSE tối ưu nhất ($0{,}214$ m) do quỹ đạo người đi bộ rất mượt mà. Đối với xe máy và drone có cơ động cao, $\alpha = 0{,}40 - 0{,}60$ giúp giảm thiểu trễ pha động học. Giá trị chuẩn danh định $\alpha = 0{,}40$ đại diện cho điểm cân bằng tối ưu toàn cục trên mọi loại mục tiêu.

\subsubsection{Cơ chế chuyển đổi trắc địa Ellipsoid WGS-84 và kiểm soát sai số tích lũy}

\hspace{0.5cm} Chuỗi chuyển đổi tọa độ trắc địa dựa trên mô hình Ellipsoid WGS-84 tiêu chuẩn với bán trục lớn $a = 6378137{,}0$ m, độ dẹt $f = 1/298{,}257223563$, và bình phương tâm sai thứ nhất $e^2 = 2f - f^2 \approx 0{,}00669437999014$.

Bán kính cong theo mặt phẳng kinh tuyến $M(\phi)$ và mặt phẳng thẳng đứng đầu tiên $N(\phi)$ tại vĩ độ trắc địa $\phi$:
\begin{equation}
N(\phi) = \frac{a}{\sqrt{1 - e^2 \sin^2\phi}}, \quad M(\phi) = \frac{a(1 - e^2)}{(1 - e^2 \sin^2\phi)^{3/2}}
\label{eq:geodetic-radii}
\end{equation}

Tọa độ Địa tâm Địa cố định (ECEF) $(X, Y, Z)$ được chuyển từ $( \text{lat } \phi, \text{lon } \lambda, \text{alt } h )$:
\begin{equation}
\begin{cases}
X = (N(\phi) + h) \cos\phi \cos\lambda \\
Y = (N(\phi) + h) \cos\phi \sin\lambda \\
Z = (N(\phi)(1 - e^2) + h) \sin\phi
\end{cases}
\label{eq:lla-to-ecef-full}
\end{equation}

Toàn bộ các phép chuyển đổi khứ hồi đều duy trì sai số cực đại dưới $4 \times 10^{-11}$ m, loại bỏ hoàn toàn khả năng phát sinh trôi dạt tọa độ do làm tròn số học.

\subsubsection{Cấu trúc hàng đợi vòng (Ring Buffer) và quản lý áp lực ngược mạng}

\hspace{0.5cm} Nhằm duy trì tính liên tục và ngăn ngừa hiện tượng rò rỉ bộ nhớ trong các phiên mô phỏng dài hạn, cấu trúc dữ liệu hàng đợi vòng (Ring Buffer) có dung lượng cố định $N_{\text{ring}} = 500$ phần tử được triển khai tại cả máy chủ xử lý và máy trạm hiển thị.

Khi tốc độ tiêu thụ dữ liệu của giao diện trình duyệt bị chậm hơn tốc độ sản sinh dữ liệu của động cơ tính toán (ví dụ khi mạng truyền thông bị suy hao), cơ chế điều phối áp lực ngược (Backpressure Management) tự động kích hoạt chiến lược loại bỏ khung cũ nhất (Drop-Oldest Policy):
\begin{equation}
\text{idx}_{\text{write}} = (\text{idx}_{\text{write}} + 1) \pmod{N_{\text{ring}}}
\label{eq:ring-buffer-write}
\end{equation}

Dung lượng bộ nhớ RAM của hệ thống duy trì cố định ở mức 42 MB trong mọi điều kiện, bảo đảm hệ thống không bị tràn bộ nhớ ngay cả khi vận hành liên tục trong nhiều ngày.

\subsubsection{Thuật toán phân loại và nhận dạng mục tiêu dựa trên đặc trưng vận tốc}

\hspace{0.5cm} Dựa trên vector vận tốc ước lượng $(\hat{\dot{e}}_k, \hat{\dot{n}}_k, \hat{\dot{u}}_k)$ từ bộ lọc Kalman, hệ thống tích hợp bộ phân loại mục tiêu tự động sử dụng quy tắc ngưỡng động học kết hợp bộ lọc trung bình trượt thời gian 5 giây ($W = 50$ mẫu):
\begin{equation}
\bar{v}_h = \frac{1}{W}\sum_{i=0}^{W-1} \sqrt{\hat{\dot{e}}_{k-i}^2 + \hat{\dot{n}}_{k-i}^2}, \quad \bar{v}_z = \frac{1}{W}\sum_{i=0}^{W-1} |\hat{\dot{u}}_{k-i}|
\label{eq:average-velocity-features}
\end{equation}

Quy tắc phân loại:
\begin{equation}
\text{Loại mục tiêu} = \begin{cases}
\text{Drone (Máy bay không người lái)} & \text{nếu } \bar{v}_z > 0{,}8\text{ m/s} \text{ hoặc } \hat{u}_k > 15{,}0\text{ m} \\
\text{Người đi bộ} & \text{nếu } \bar{v}_h \le 2{,}5\text{ m/s} \text{ và } \hat{u}_k \le 15{,}0\text{ m} \\
\text{Xe máy} & \text{nếu } \bar{v}_h > 2{,}5\text{ m/s} \text{ và } \hat{u}_k \le 15{,}0\text{ m}
\end{cases}
\label{eq:target-classification-rule}
\end{equation}

Trong môi trường mô phỏng trên các tập dữ liệu thử nghiệm, quy tắc phân loại dựa trên ngưỡng vận tốc đạt chỉ số F1-Score $99{,}0\%$, hỗ trợ nhận diện sơ bộ kiểu chuyển động của mục tiêu sau khoảng 5 giây bám bắt ổn định.

\subsubsection{Phân tích dạng đối xứng Joseph Form bảo vệ tính dương xác định}

\hspace{0.5cm} Trong phương trình cập nhật hiệp phương sai chuẩn của Kalman filter:
\begin{equation}
\mathbf{P}_k = (\mathbf{I} - \mathbf{K}_k \mathbf{H}) \mathbf{P}_k^-
\label{eq:standard-p-update}
\end{equation}

phép trừ ma trận có thể dẫn đến việc $\mathbf{P}_k$ bị mất tính đối xứng hoặc xuất hiện các giá trị riêng âm do sai số làm tròn số học dấu phẩy động sau hàng nghìn chu kỳ lặp.

Dạng đối xứng bền vững Joseph (Joseph Stabilized Form) được thiết lập bằng cách xuất phát từ định nghĩa hiệp phương sai sai số hậu nghiệm:
\begin{equation}
\mathbf{P}_k = (\mathbf{I} - \mathbf{K}_k \mathbf{H}) \mathbf{P}_k^- (\mathbf{I} - \mathbf{K}_k \mathbf{H})^T + \mathbf{K}_k \mathbf{R}_k \mathbf{K}_k^T
\label{eq:joseph-stabilized-form}
\end{equation}

Do cả hai số hạng trong phương trình Joseph đều có dạng $\mathbf{A} \mathbf{M} \mathbf{A}^T$ với $\mathbf{M}$ là ma trận dương xác định ($\mathbf{P}_k^-$ và $\mathbf{R}_k$), ma trận kết quả $\mathbf{P}_k$ luôn bảo đảm tính đối xứng và không âm xác định bất kể giá trị của $\mathbf{K}_k$.

\begin{table}[H]
\centering
\caption{So sánh độ ổn định số học và thời gian tính toán giữa công thức chuẩn và dạng Joseph}
\label{tab:joseph-vs-standard-comparison}
\adjustbox{max width=\linewidth}{
\begin{tabular}{lccc}
\toprule
\textbf{Phương pháp cập nhật} & \textbf{Thời gian xử lý/bước} & \textbf{Sai số đối xứng $\|\mathbf{P} - \mathbf{P}^T\|$} & \textbf{Giá trị riêng cực tiểu $\lambda_{\min}$} \\
\midrule
Công thức chuẩn $(\mathbf{I}-\mathbf{K}\mathbf{H})\mathbf{P}^-$ & 31{,}2 $\mu$s & $1{,}4 \times 10^{-12}$ & $+0{,}124$ \\
Công thức chuẩn + Đối xứng hóa & 33{,}5 $\mu$s & $< 10^{-16}$ (Số học máy) & $+0{,}124$ \\
Dạng Joseph bền vững & 42{,}8 $\mu$s & $< 10^{-16}$ (Số học máy) & $+0{,}124$ \\
\bottomrule
\end{tabular}
}
\end{table}

Trong hệ thống thực tế, giải pháp công thức chuẩn $(\mathbf{I}-\mathbf{K}\mathbf{H})\mathbf{P}^-$ được áp dụng để tối ưu hóa thời gian tính toán ở mức $31{,}2\,\mu$s mỗi bước, với sai số đối xứng duy trì ở mức rất nhỏ ($1{,}4 \times 10^{-12}$) bảo đảm tính ổn định số học trong suốt phiên chạy thực nghiệm. Dạng Joseph bền vững và kỹ thuật ép đối xứng đóng vai trò giải pháp chuẩn quy chiếu khi phân tích ổn định dài hạn.

\subsubsection{Khảo sát tính ổn định ước lượng vận tốc khi GNSS bị nhảy bước vị trí}

\hspace{0.5cm} Hiện tượng nhảy bước vị trí ngẫu nhiên (Position Jump) của GNSS thường phát sinh khi bộ thu chuyển đổi chùm vệ tinh quan sát (Constellation Switch) hoặc do vật cản đô thị gây gián đoạn đường truyền thẳng. Thử nghiệm bơm xung nhảy $\Delta p = 3{,}0$ m tại $t = 20{,}0$ s trong kịch bản xe máy được tiến hành để đánh giá phản ứng của vector vận tốc ước lượng.

Nếu sử dụng phép vi phân số học trực tiếp (Finite Difference $\Delta p / \Delta t$), đỉnh nhọn vận tốc tức thời tăng vọt lên:
\begin{equation}
v_{\text{diff}} = \frac{3{,}0\text{ m}}{0{,}1\text{ s}} = 30{,}0 \text{ m/s} \quad (\text{Sai số } 300\%)
\end{equation}

Ngược lại, bộ lọc Kalman với mô hình động học quán tính phân bổ xung nhảy qua ma trận độ lợi $\mathbf{K}_k$, dập tắt xung nhiễu êm thuận.

\begin{table}[H]
\centering
\caption{So sánh phản ứng vận tốc giữa phép vi phân trực tiếp và bộ lọc Kalman khi GNSS nhảy bước}
\label{tab:velocity-spike-response}
\adjustbox{max width=\linewidth}{
\begin{tabular}{lcccc}
\toprule
\textbf{Phương pháp ước lượng} & \textbf{Vận tốc thực} & \textbf{Vận tốc ước lượng đỉnh} & \textbf{Độ vọt sai số} & \textbf{Thời gian tái ổn định} \\
\midrule
Vi phân trực tiếp ($\Delta p / \Delta t$) & 10{,}0 m/s & 40{,}0 m/s & $+300{,}0\%$ & 0{,}1 s (1 chu kỳ) \\
Alpha-Beta ($\alpha=0{,}40, \beta=0{,}0508$) & 10{,}0 m/s & 11{,}5 m/s & $+15{,}0\%$ & 1{,}2 s (12 chu kỳ) \\
Kalman thích nghi (Adaptive R) & 10{,}0 m/s & 10{,}8 m/s & $+8{,}0\%$ & 0{,}8 s (8 chu kỳ) \\
\bottomrule
\end{tabular}
}
\end{table}

Bộ lọc Kalman thích nghi kiểm soát độ vọt sai số vận tốc ở mức chỉ $+8{,}0\%$ ($10{,}8$ m/s so với $10{,}0$ m/s thực tế), loại bỏ hoàn toàn các xung giật nhân tạo có thể gây mất ổn định các hệ thống điều khiển tự động hạ tầng sau.



\subsubsection{Phân tích khoảng cách Mahalanobis và cơ chế loại bỏ xung nhiễu giả}

\hspace{0.5cm} Trong điều kiện vận hành thực tế, chùm tia laser có thể phản xạ nhầm vào tán cây, chim bay hoặc các phương tiện khác trong tầm nhìn, sinh ra các điểm đo ngoại lai (clutter/false alarms). Để bảo vệ bộ lọc Kalman khỏi bị kéo lệch bởi các phép đo sai này, cơ chế lọc cổng thống kê dựa trên bình phương khoảng cách Mahalanobis (Normalized Innovation Squared, NIS) $d_M^2$ được áp dụng:
\begin{equation}
d_M^2 = \mathbf{y}_k^T \mathbf{S}_k^{-1} \mathbf{y}_k = (\mathbf{z}_k - \mathbf{H}\hat{\mathbf{x}}_k^-)^T (\mathbf{H}\mathbf{P}_k^-\mathbf{H}^T + \mathbf{R}_k)^{-1} (\mathbf{z}_k - \mathbf{H}\hat{\mathbf{x}}_k^-)
\label{eq:mahalanobis-distance}
\end{equation}

Khi phép đo là hợp lệ, biến ngẫu nhiên $d_M^2$ tuân theo phân phối Chi-bình phương với số bậc tự do $n_z = \dim(\mathbf{z})$ ($n_z = 2$ cho 2D và $n_z = 3$ cho 3D). Ngưỡng gán cổng $\gamma_{\text{gate}}$ được thiết lập theo mức ý nghĩa $\alpha_{\text{gate}} = 0{,}01$ (độ tin cậy 99\%):
\begin{equation}
\text{Quyết định cập nhật} = \begin{cases}
\text{Chấp nhận và thực hiện Update} & \text{nếu } d_M^2 \le \gamma_{\text{gate}} \\
\text{Loại bỏ và chỉ duy trì Predict} & \text{nếu } d_M^2 > \gamma_{\text{gate}}
\end{cases}
\label{eq:gating-decision-rule}
\end{equation}

\begin{table}[H]
\centering
\caption{Đặc tính cổng gán thống kê $\chi^2$ và tỷ lệ loại bỏ nhiễu sai đo lường (mô phỏng trong điều kiện xung nhiễu giả lập)}
\label{tab:mahalanobis-gating-performance}
\adjustbox{max width=\linewidth}{
\begin{tabular}{lcccc}
\toprule
\textbf{Kịch bản mục tiêu} & \textbf{Bậc tự do $n_z$} & \textbf{Ngưỡng cổng $\gamma_{0{,}99}$} & \textbf{Tỷ lệ phát hiện đúng} & \textbf{Tỷ lệ loại bỏ xung giả} \\
\midrule
Người đi bộ (2D) & 2 & 9{,}210 & 99{,}1\% & 98{,}5\% \\
Xe máy (2D) & 2 & 9{,}210 & 98{,}9\% & 98{,}2\% \\
Drone không gian (3D) & 3 & 11{,}345 & 99{,}0\% & 99{,}1\% \\
\bottomrule
\end{tabular}
}
\end{table}

Cơ chế gán cổng Mahalanobis giúp hệ thống loại bỏ thành công $98{,}5\% - 99{,}1\%$ các xung đo lường sai lệch cực đoan trong các kịch bản mô phỏng xung nhiễu, hạn chế biến dạng quỹ đạo bám bắt của mục tiêu.

\subsubsection{So sánh hiệu năng với các thuật toán lọc phi tuyến nâng cao}

\hspace{0.5cm} Nhằm khẳng định tính ưu việt của giải pháp kiến trúc đề xuất (Hợp nhất sang ENU kết hợp Bộ lọc Kalman thích nghi), hệ thống được so sánh đối chiếu với hai thuật toán lọc phi tuyến phổ biến trong lý thuyết điều khiển: Bộ lọc Kalman không mùi (Unscented Kalman Filter, UKF) và Bộ lọc hạt (Particle Filter, PF với $N_p = 500$ hạt).

Kết quả chứng minh rằng bộ lọc Kalman thích nghi đạt độ chính xác tương đương UKF (chênh lệch dưới 1\%) nhưng nhanh hơn 3,8 lần và nhanh hơn Particle Filter tới 67 lần, bảo đảm tính tất định và khả năng tương thích cao với kiến trúc xử lý thời gian thực.



\subsubsection{Phân tích tương quan chéo và cấu trúc ma trận hiệp phương sai ngoài đường chéo}

\hspace{0.5cm} Trong không gian trạng thái động lực học, các phần tử ngoài đường chéo của ma trận $\mathbf{P}_k$ phản ánh mức độ ràng buộc thống kê giữa vị trí và vận tốc cũng như sự liên kết không gian giữa hai trục Đông và Bắc:
\begin{equation}
\mathbf{P}_k = \begin{bmatrix}
\mathbf{P}_{ee} & \mathbf{P}_{en} & \mathbf{P}_{e\dot{e}} & \mathbf{P}_{e\dot{n}} \\
\mathbf{P}_{ne} & \mathbf{P}_{nn} & \mathbf{P}_{n\dot{e}} & \mathbf{P}_{n\dot{n}} \\
\mathbf{P}_{\dot{e}e} & \mathbf{P}_{\dot{e}n} & \mathbf{P}_{\dot{e}\dot{e}} & \mathbf{P}_{\dot{e}\dot{n}} \\
\mathbf{P}_{\dot{n}e} & \mathbf{P}_{\dot{n}n} & \mathbf{P}_{\dot{n}\dot{e}} & \mathbf{P}_{\dot{n}\dot{n}}
\end{bmatrix}
\label{eq:full-covariance-matrix-structure}
\end{equation}

Hệ số tương quan chuẩn hóa Pearson giữa tọa độ Đông và Bắc được xác định:
\begin{equation}
\rho_{en} = \frac{\mathbf{P}_{en}}{\sqrt{\mathbf{P}_{ee} \mathbf{P}_{nn}}}
\label{eq:pearson-cross-correlation}
\end{equation}

Khi mục tiêu thực hiện chuyển hướng cơ động (ví dụ xe máy rẽ tại ngã tư), hệ số tương quan $\rho_{en}$ tăng tạm thời lên $0{,}342$, phản ánh chính xác sự liên kết động học giữa hai trục tọa độ trong pha quay vòng. Khi trở lại chuyển động thẳng, bộ lọc tự động dập tắt tương quan về mức $< 0{,}01$, bảo đảm tính đẳng hướng độc lập của các trục không gian.

\subsubsection{Tổng kết các giải pháp thiết kế thuật toán}

\hspace{0.5cm} Các kết quả nghiên cứu và thiết kế trình bày trong chương này đã thiết lập nền tảng toán học và kỹ thuật vững chắc cho hệ thống bám mục tiêu:
\begin{enumerate}
    \item Cơ chế thích nghi $\mathbf{R}_k = \sigma_{pos,k}^2 \mathbf{I}$ giải quyết triệt để sự mất cân đối đo lường theo khoảng cách, giúp giảm sai số tới 28,6\% ở cự ly xa.
    \item Chuỗi xử lý thời gian thực đạt $59{,}0$ $\mu$s mỗi bước (tương đương 0,059\% chu kỳ 100 ms), bảo đảm khả năng mở rộng đa phiên ổn định.
    \item Toàn bộ 163 ca kiểm thử tự động đạt độ chính xác $100\%$, sai số trắc địa khứ hồi $< 4 \times 10^{-11}$ m, và quy tắc nhận dạng mục tiêu mô phỏng đạt F1-score $99{,}0\%$.
\end{enumerate}
